\documentclass[12pt]{article} % Тип документа (стаття) та розмір шрифту (12pt)
\usepackage[T2A]{fontenc}      % Кодування шрифтів для відображення кирилиці
\usepackage[utf8]{inputenc}    % Кодування вхідного файлу (стандарт для укр.мови)
\usepackage[ukrainian]{babel}  % Підключення словників та правил переносу укр.мови
\usepackage{amsmath, amssymb}  % Пакети для математичних формул та символів
\usepackage{array}             % Пакет для покращеної роботи з таблицями

\title{Практична робота: Математичні формули} 
\author{Яцюк Вікторія}                          
\date{\today}                                

\begin{document}
\maketitle % Команда, яка фактично малює заголовок, автора та дату на сторінці 
 \section* {Квадратне рівняння}  

Квадратне рівняння має вигляд:

\begin{equation}
ax^2 + bx + c = 0, \quad a \neg 0
\end{equation}

Його розв'язки залежать від значення дискримінанта $D$. Нижче наведено узагальнюючу таблицю:



\end{document} % Кінець документа
